\documentclass{amsart}
% For dealing with references we use the comment environment
\usepackage{verbatim}
\newenvironment{reference}{\comment}{\endcomment}
%\newenvironment{reference}{}{}
\newenvironment{slogan}{\comment}{\endcomment}
\newenvironment{history}{\comment}{\endcomment}

% For commutative diagrams we use Xy-pic
\usepackage[all]{xy}

% We use 2cell for 2-commutative diagrams.
\xyoption{2cell}
\UseAllTwocells

% We use multicol for the list of chapters between chapters
\usepackage{multicol}

% This is generally recommended for better output
\usepackage{lmodern}
\usepackage[T1]{fontenc}

% For cross-file-references

% Package for hypertext links:
\usepackage{hyperref}

% For any local file, say "hello.tex" you want to link to please

% Theorem environments.
%
\theoremstyle{plain}
\newtheorem{theorem}[subsection]{Theorem}
\newtheorem{proposition}[subsection]{Proposition}
\newtheorem{lemma}[subsection]{Lemma}

\theoremstyle{definition}
\newtheorem{definition}[subsection]{Definition}
\newtheorem{example}[subsection]{Example}
\newtheorem{exercise}[subsection]{Exercise}
\newtheorem{situation}[subsection]{Situation}

\theoremstyle{remark}
\newtheorem{remark}[subsection]{Remark}
\newtheorem{remarks}[subsection]{Remarks}

\numberwithin{equation}{subsection}

% Macros
%
\def\lim{\mathop{\mathrm{lim}}\nolimits}
\def\colim{\mathop{\mathrm{colim}}\nolimits}
\def\Spec{\mathop{\mathrm{Spec}}}
\def\Hom{\mathop{\mathrm{Hom}}\nolimits}
\def\Ext{\mathop{\mathrm{Ext}}\nolimits}
\def\SheafHom{\mathop{\mathcal{H}\!\mathit{om}}\nolimits}
\def\SheafExt{\mathop{\mathcal{E}\!\mathit{xt}}\nolimits}
\def\Sch{\mathit{Sch}}
\def\Mor{\mathop{\mathrm{Mor}}\nolimits}
\def\Ob{\mathop{\mathrm{Ob}}\nolimits}
\def\Sh{\mathop{\mathit{Sh}}\nolimits}
\def\NL{\mathop{N\!L}\nolimits}
\def\CH{\mathop{\mathrm{CH}}\nolimits}
\def\proetale{{pro\text{-}\acute{e}tale}}
\def\etale{{\acute{e}tale}}
\def\QCoh{\mathit{QCoh}}
\def\Ker{\mathop{\mathrm{Ker}}}
\def\Im{\mathop{\mathrm{Im}}}
\def\Coker{\mathop{\mathrm{Coker}}}
\def\Coim{\mathop{\mathrm{Coim}}}

% Boxtimes
%
\DeclareMathSymbol{\boxtimes}{\mathbin}{AMSa}{"02}

%
% Macros for moduli stacks/spaces
%
\def\QCohstack{\mathcal{QC}\!\mathit{oh}}
\def\Cohstack{\mathcal{C}\!\mathit{oh}}
\def\Spacesstack{\mathcal{S}\!\mathit{paces}}
\def\Quotfunctor{\mathrm{Quot}}
\def\Hilbfunctor{\mathrm{Hilb}}
\def\Curvesstack{\mathcal{C}\!\mathit{urves}}
\def\Polarizedstack{\mathcal{P}\!\mathit{olarized}}
\def\Complexesstack{\mathcal{C}\!\mathit{omplexes}}
% \Pic is the operator that assigns to X its picard group, usage \Pic(X)
% \Picardstack_{X/B} denotes the Picard stack of X over B
% \Picardfunctor_{X/B} denotes the Picard functor of X over B
\def\Pic{\mathop{\mathrm{Pic}}\nolimits}
\def\Picardstack{\mathcal{P}\!\mathit{ic}}
\def\Picardfunctor{\mathrm{Pic}}
\def\Deformationcategory{\mathcal{D}\!\mathit{ef}}

\setlength{\emergencystretch}{3em}
\begin{document}
\title[Stacks Project, Tag 0BST]{299 --- Etaleness of inertia invariants for finite group actions on arbitrary rings}
\maketitle
\noindent Copyright (C) 2005--2025 Johan de Jong. Stacks Project authors. Source: \href{https://stacks.math.columbia.edu/tag/0BST}{Tag 0BST}.

\noindent Editorial difficulty note (not author text). Difficulty H2: The proof must show the invariant-ring map is etale at a selected prime without normality or finite-type assumptions. Etale localization splits the integral fibre, and a further equivariant clopen refinement makes the group permute factors; the inertia subgroup then identifies a factor of the full invariant ring..

\noindent Editorial provenance note (not author text). History: excerpt from revision \texttt{a0444\allowbreak 6e57e\allowbreak c1fbc\allowbreak 252a8\allowbreak 71afc\allowbreak ec775\allowbreak 2fb28\allowbreak 07b14}, pione.tex, lines 2664--2778. Reference presentation was adapted; original-author context and core support blocks were retained or added. The mathematical wording is unchanged. Licensed under GNU FDL 1.2 or later; no invariant sections or cover texts. See ../licenses/COPYING. Context blocks reproduce author definitions and supporting results; they are not additional counted problems.

\begin{lemma}
\label{lemma-inertia-invariants-etale}
Let $G$ be a finite group acting on a ring $R$. Let $\mathfrak q \subset R$
be a prime. Set
$$
I = \{\sigma \in G \mid \sigma(\mathfrak q) = \mathfrak q
\text{ and } \sigma \bmod \mathfrak q = \text{id}_\mathfrak q\}
$$
Then $R^G \to R^I$ is \'etale at $R^I \cap \mathfrak q$.
\end{lemma}

\begin{proof}
The strategy of the proof is to use \'etale localization to
reduce to the case where $R \to R^I$ is a local isomorphism at
$R^I \cap \mathfrak p$.
Let $R^G \to A$ be an \'etale ring map. We claim that if the result
holds for the action of $G$ on $A \otimes_{R^G} R$ and some prime
$\mathfrak q'$ of $A \otimes_{R^G} R$ lying over $\mathfrak q$, then
the result is true.

\medskip\noindent
To check this, note that since $R^G \to A$ is flat we have
$A = (A \otimes_{R^G} R)^G$, see More on Algebra,
Lemma \href{https://stacks.math.columbia.edu/tag/0BRH}{lemma base change invariants (Tag 0BRH)}.
By Lemma \href{https://stacks.math.columbia.edu/tag/0BSS}{lemma inertia base change (Tag 0BSS)} the group $I$ does not change.
Then a second application of More on Algebra,
Lemma \href{https://stacks.math.columbia.edu/tag/0BRH}{lemma base change invariants (Tag 0BRH)}
shows that $A \otimes_{R^G} R^I = (A \otimes_{R^G} R)^I$
(because $R^I \to A \otimes_{R^G} R^I$ is flat).
Thus
$$
\xymatrix{
\Spec((A \otimes_{R^G} R)^I) \ar[d] \ar[r] & \Spec(R^I) \ar[d] \\
\Spec(A) \ar[r] & \Spec(R^G)
}
$$
is cartesian and the horizontal arrows are \'etale. Thus if the
left vertical arrow is \'etale in some open neighbourhood $W$ of
$(A \otimes_{R^G} R)^I \cap \mathfrak q'$, then the right vertical
arrow is \'etale at the points of the (open) image of $W$ in
$\Spec(R^I)$, see
Descent, Lemma \href{https://stacks.math.columbia.edu/tag/05B5}{lemma smooth permanence (Tag 05B5)}. In particular
the morphism $\Spec(R^I) \to \Spec(R^G)$ is \'etale at $R^I \cap \mathfrak q$.

\medskip\noindent
Let $\mathfrak p = R^G \cap \mathfrak q$.
By More on Algebra, Lemma \href{https://stacks.math.columbia.edu/tag/0BRI}{lemma one orbit (Tag 0BRI)}
the fibre of $\Spec(R) \to \Spec(R^G)$ over $\mathfrak p$ is
finite. Moreover the residue field extensions at these points
are algebraic, normal, with finite automorphism groups by
More on Algebra, Lemma \href{https://stacks.math.columbia.edu/tag/0BRJ}{lemma one orbit geometric (Tag 0BRJ)}.
Thus we may apply
More on Morphisms,
Lemma \href{https://stacks.math.columbia.edu/tag/0BSR}{lemma etale makes integral split (Tag 0BSR)}
to the integral ring map $R^G \to R$ and the prime $\mathfrak p$.
Combined with the claim above we reduce to the case where
$R = A_1 \times \ldots \times A_n$ with each $A_i$ having a single
prime $\mathfrak q_i$ lying over $\mathfrak p$ such that the
residue field extensions $\kappa(\mathfrak q_i)/\kappa(\mathfrak p)$
are purely inseparable. Of course $\mathfrak q$ is one of
these primes, say $\mathfrak q = \mathfrak q_1$.

\medskip\noindent
It may not be the case that $G$ permutes the factors $A_i$
(this would be true if the spectrum of $A_i$ were connected,
for example if $R^G$ was local). This we can fix as follows;
we suggest the reader think this through for themselves, perhaps
using idempotents instead of topology.
Recall that the product decomposition gives a corresponding
disjoint union decomposition of $\Spec(R)$ by open and closed
subsets $U_i$. Since $G$ is finite, we can refine this covering
by a finite disjoint union decomposition
$\Spec(R) = \coprod_{j \in J} W_j$ by open
and closed subsets $W_j$, such that for all $j \in J$ there exists
a $j' \in J$ with $\sigma(W_j) = W_{j'}$. The union of the
$W_j$ not meeting $\{\mathfrak q_1, \ldots, \mathfrak q_n\}$
is a closed subset not meeting the fibre over $\mathfrak p$
hence maps to a closed subset of $\Spec(R^G)$ not meeting
$\mathfrak p$ as $\Spec(R) \to \Spec(R^G)$ is closed.
Hence after replacing $R^G$ by a principal localization
(permissible by the claim) we may assume each $W_j$ meets
one of the points $\mathfrak q_i$. Then we set $U_i = W_j$
if $\mathfrak q_i \in W_j$. The corresponding product decomposition
$R = A_1 \times \ldots \times A_n$ is one
where $G$ permutes the factors $A_i$.

\medskip\noindent
Thus we may assume we have a product decomposition
$R = A_1 \times \ldots \times A_n$ compatible with $G$-action,
where each $A_i$ has a single prime $\mathfrak q_i$ lying
over $\mathfrak p$ and the field extensions
$\kappa(\mathfrak q_i)/\kappa(\mathfrak p)$ are purely inseparable.
Write $A' = A_2 \times \ldots \times A_n$ so that
$$
R = A_1 \times A'
$$
Since $\mathfrak q = \mathfrak q_1$ we find that every
$\sigma \in I$ preserves the product decomposition above.
Hence
$$
R^I = (A_1)^I \times (A')^I
$$
Observe that $I = D = \{\sigma \in G \mid \sigma(\mathfrak q) = \mathfrak q\}$
because $\kappa(\mathfrak q)/\kappa(\mathfrak p)$ is purely inseparable.
Since the action of $G$ on primes over $\mathfrak p$ is transitive
(More on Algebra, Lemma \href{https://stacks.math.columbia.edu/tag/0BRI}{lemma one orbit (Tag 0BRI)})
we conclude that, the index of $I$ in $G$ is $n$ and we can write
$G = eI \amalg \sigma_2I \amalg \ldots \amalg \sigma_nI$ so that
$A_i = \sigma_i(A_1)$ for $i = 2, \ldots, n$. It follows that
$$
R^G = (A_1)^I.
$$
Thus the map $R^G \to R^I$ is \'etale at $R^I \cap \mathfrak q$
and the proof is complete.
\end{proof}
\end{document}
